% Preámbulo del informe técnico de Alicanto.
% Compila con pdfLaTeX (también en Overleaf, con el compilador por defecto). No necesita shell-escape.

\usepackage[T1]{fontenc}
\usepackage[spanish,es-tabla,es-noquoting,es-noshorthands,es-nolists,es-noindentfirst,es-nosectiondot]{babel}
\usepackage[top=2.7cm,bottom=3cm,left=2.9cm,right=2.9cm,headsep=0.7cm,footskip=1.2cm]{geometry}

% --- Tipografía -------------------------------------------------------------
\usepackage[scaled=.98]{XCharter}
\usepackage[scaled=0.95]{sourcesanspro}
\usepackage[scaled=0.86,varqu,varl]{zi4}
\usepackage{amsmath}
\usepackage[charter,vvarbb,scaled=1.02]{newtxmath}
\usepackage[final,tracking=smallcaps]{microtype}
\linespread{1.06}

% --- Color ------------------------------------------------------------------
\usepackage[dvipsnames,table]{xcolor}
\definecolor{tinta}{HTML}{17140F}      % texto
\definecolor{oro}{HTML}{C8941F}        % acento (el oro de la demo, oscurecido para papel blanco)
\definecolor{oroclaro}{HTML}{E8B23A}   % el oro de la demo, para rellenos
\definecolor{oroscuro}{HTML}{8A6410}   % enlaces
\definecolor{ceniza}{HTML}{6B665C}     % texto secundario
\definecolor{cenizaclara}{HTML}{BDB8AD}
\definecolor{humo}{HTML}{F4F1EA}       % fondos
\definecolor{noche}{HTML}{0D0C0A}      % portada
\definecolor{marfil}{HTML}{F3EFE4}
\definecolor{rojoeq}{HTML}{D44E40}
\definecolor{moradoeq}{HTML}{9267CF}
\definecolor{azuleq}{HTML}{4B81DA}
\color{tinta}

% --- Tablas, listas, unidades ----------------------------------------------
\usepackage{array,booktabs,tabularx,longtable,multirow}
% Las filas generadas por hacer_datos.py se leen aquí, una vez, y se guardan como macros.
% Dentro de una tabla no se puede leer un archivo (\input falla justo antes de \bottomrule).
\usepackage{catchfile}
\newcommand{\cargartabla}[1]{\expandafter\CatchFileDef\csname datos@#1\endcsname{datos/#1.tex}{}}
\newcommand{\tabladatos}[1]{\csname datos@#1\endcsname}
\cargartabla{tab_grafo}\cargartabla{tab_ell}\cargartabla{tab_busqueda_rot}\cargartabla{tab_busqueda_bits}
\cargartabla{tab_cambios}\cargartabla{tab_sistemas}\cargartabla{tab_mecanismo}\cargartabla{tab_normales}
\cargartabla{tab_sistemas_familia}\cargartabla{tab_adaptativo}\cargartabla{tab_adaptativo_familia}
\cargartabla{tab_fragmentos_matriz}\cargartabla{tab_modelo}\cargartabla{tab_fragmentos}
\usepackage{enumitem}
\setlist{itemsep=0.25ex,topsep=0.6ex,parsep=0.3ex,leftmargin=1.5em}
\setlist[description]{font=\normalfont\bfseries,leftmargin=1.5em,style=nextline}
\usepackage{siunitx}
\sisetup{output-decimal-marker={,},group-separator={.},group-minimum-digits=4,group-digits=integer,
         detect-weight=true,detect-family=true,mode=text,range-phrase={~a~}}
\renewcommand{\arraystretch}{1.18}
\newcolumntype{L}{>{\raggedright\arraybackslash}X}
\newcolumntype{P}[1]{>{\raggedright\arraybackslash}p{#1}}
\newcolumntype{R}[1]{>{\raggedleft\arraybackslash}p{#1}}

% --- Figuras ----------------------------------------------------------------
\usepackage{graphicx}
\usepackage{tikz}
\usetikzlibrary{arrows.meta,positioning,calc,fit,backgrounds,shapes.geometric,shapes.misc,
                decorations.pathreplacing,matrix,babel}
\usepackage{pgfplots}
\usepgfplotslibrary{groupplots}
\pgfplotsset{compat=1.18,
  every axis/.append style={
    axis line style={ceniza!70}, tick style={ceniza!70}, tick align=outside,
    label style={font=\small}, tick label style={font=\footnotesize},
    title style={font=\small\sffamily\bfseries,yshift=-2pt},
    legend style={draw=none,fill=none,font=\footnotesize,legend cell align=left},
    /pgf/number format/use comma, /pgf/number format/1000 sep={.},
  }}
\tikzset{
  >={Stealth[length=5pt]},
  caja/.style={draw=tinta,line width=0.6pt,rounded corners=2pt,fill=white,align=center,
               inner sep=5pt,font=\small\sffamily},
  cajaoro/.style={caja,draw=oro,fill=oroclaro!18},
  cajagris/.style={caja,draw=ceniza!70,fill=humo},
  flecha/.style={->,line width=0.6pt,draw=tinta},
  flechagris/.style={->,line width=0.6pt,draw=ceniza},
  nota/.style={font=\footnotesize,text=ceniza,align=left},
  decision/.style={draw=tinta,line width=0.6pt,diamond,aspect=2.4,fill=white,align=center,
                   inner sep=1.5pt,font=\footnotesize\sffamily},
  casilla frenado/.style={fill=tinta,draw=tinta,line width=0.3pt},
  casilla escapa/.style={fill=oroclaro,draw=oro,line width=0.3pt},
  casilla rotada/.style={fill=cenizaclara,draw=ceniza,line width=0.3pt},
  casilla vacia/.style={fill=white,draw=ceniza!60,line width=0.3pt},
}

% --- Código y comandos -------------------------------------------------------
\usepackage{fvextra}
\DefineVerbatimEnvironment{consola}{Verbatim}{fontsize=\small,breaklines=true,breakanywhere=true,
  frame=leftline,framerule=1.4pt,framesep=8pt,rulecolor=\color{cenizaclara},xleftmargin=10pt,
  breaksymbolleft={}}
\usepackage[ruled,vlined,linesnumbered,spanish,onelanguage]{algorithm2e}
\SetAlFnt{\small}
\SetAlCapFnt{\small\sffamily}
\SetAlCapNameFnt{\small}
\SetKwComment{Com}{$\triangleright$\ }{}
\SetCommentSty{textit}
\DontPrintSemicolon

% --- Recuadros ----------------------------------------------------------------
\usepackage[most]{tcolorbox}
\newtcolorbox{clave}[1][]{enhanced,breakable,frame hidden,boxrule=0pt,arc=0pt,
  colback=oroclaro!13,borderline west={2.2pt}{0pt}{oro},left=10pt,right=8pt,top=6pt,bottom=6pt,
  before skip=0.9\baselineskip,after skip=0.9\baselineskip,fontupper=\small,#1}
\newtcolorbox{limite}[1][]{enhanced,breakable,frame hidden,boxrule=0pt,arc=0pt,
  colback=humo,borderline west={2.2pt}{0pt}{ceniza},left=10pt,right=8pt,top=6pt,bottom=6pt,
  before skip=0.9\baselineskip,after skip=0.9\baselineskip,fontupper=\small,#1}

% --- Encabezados y títulos ------------------------------------------------------
\usepackage[automark,headsepline=0.4pt]{scrlayer-scrpage}
\clearpairofpagestyles
\ihead{\footnotesize\sffamily\color{ceniza}Alicanto \textperiodcentered\ informe final y explicación técnica}
\ohead{\footnotesize\sffamily\color{ceniza}\leftmark}
\cfoot{\footnotesize\sffamily\color{ceniza}\pagemark}
\setkomafont{headsepline}{\color{cenizaclara}}
\setkomafont{disposition}{\sffamily\bfseries\color{tinta}}
\setkomafont{section}{\Large}
\setkomafont{subsection}{\large}
\setkomafont{subsubsection}{\normalsize}
\setkomafont{paragraph}{\normalsize\sffamily\bfseries}
\setkomafont{caption}{\small}
\setkomafont{captionlabel}{\small\sffamily\bfseries}
\setcapindent{0pt}
\renewcommand*{\sectionformat}{\textcolor{oro}{\thesection}\enskip}
\renewcommand*{\subsectionformat}{\textcolor{oro}{\thesubsection}\enskip}
\renewcommand*{\subsubsectionformat}{\textcolor{oro}{\thesubsubsection}\enskip}
\RedeclareSectionCommand[beforeskip=1.6\baselineskip,afterskip=0.6\baselineskip]{section}
\RedeclareSectionCommand[beforeskip=1.2\baselineskip,afterskip=0.35\baselineskip]{subsection}
\RedeclareSectionCommand[beforeskip=0.9\baselineskip,afterskip=0.2\baselineskip]{subsubsection}
\setcounter{secnumdepth}{3}
\setcounter{tocdepth}{2}

% Portadilla de cada parte: una página propia, con una línea de qué trae.
\newcommand{\parte}[3]{%
  \clearpage\thispagestyle{empty}\refstepcounter{part}%
  \addcontentsline{toc}{part}{Parte \thepart. #2}%
  \vspace*{0.26\textheight}%
  {\sffamily\color{oro}\large\textls[120]{PARTE #1}\par}\vspace{0.6em}%
  {\sffamily\bfseries\fontsize{30}{34}\selectfont #2\par}\vspace{1.1em}%
  {\color{oro}\rule{3.2cm}{1.6pt}\par}\vspace{1.3em}%
  \begin{minipage}{0.78\textwidth}\raggedright\large\color{ceniza}#3\end{minipage}%
  \clearpage}

% --- Enlaces y referencias cruzadas ----------------------------------------------
\PassOptionsToPackage{obeyspaces,spaces}{url}
\usepackage{xurl}
\usepackage[colorlinks=true,linkcolor=tinta,citecolor=oroscuro,urlcolor=oroscuro,
            pdfauthor={team-47, Platanus Hack 26},
            pdftitle={Alicanto: informe final y explicación técnica},
            pdfsubject={Platanus Hack 26, track AI Security}]{hyperref}
\usepackage[spanish,nameinlink,noabbrev]{cleveref}
\crefname{algocf}{algoritmo}{algoritmos}
\Crefname{algocf}{Algoritmo}{Algoritmos}
\crefname{section}{sección}{secciones}
\Crefname{section}{Sección}{Secciones}
\crefname{table}{tabla}{tablas}
\Crefname{table}{Tabla}{Tablas}
\urlstyle{tt}

% Cuadritos de leyenda que se pueden usar dentro de un \caption.
\DeclareRobustCommand{\cuadro}[1]{\textcolor{#1}{\rule{0.62em}{0.62em}}}
\DeclareRobustCommand{\cuadrovacio}{{\setlength{\fboxsep}{0pt}\setlength{\fboxrule}{0.4pt}%
  \fcolorbox{ceniza}{white}{\rule{0pt}{0.54em}\rule{0.54em}{0pt}}}}

% --- Atajos del texto ---------------------------------------------------------------
\newcommand{\apref}[1]{apéndice~\ref{#1}}     % los apéndices se citan por su letra
\newcommand{\cod}[1]{\nolinkurl{#1}}          % código y rutas: no hay que escapar _ ni /
\newcommand{\esp}[1]{\emph{#1}}                % extranjerismos
\newcommand{\bits}[1]{\num{#1}~bits}
\DeclareMathOperator{\clip}{clip}
\DeclareMathOperator{\diag}{diag}
\newcommand{\lrot}{\lambda_{\mathrm{rot}}}
\newcommand{\brot}{b_{\mathrm{rot}}}

% Caracteres Unicode que aparecen en textos copiados del repositorio.
\DeclareUnicodeCharacter{2212}{\ensuremath{-}}
\DeclareUnicodeCharacter{2192}{\ensuremath{\to}}
\DeclareUnicodeCharacter{2265}{\ensuremath{\geq}}
\DeclareUnicodeCharacter{2264}{\ensuremath{\leq}}
\DeclareUnicodeCharacter{00D7}{\ensuremath{\times}}
\DeclareUnicodeCharacter{2032}{\ensuremath{'}}
\DeclareUnicodeCharacter{00B7}{\ensuremath{\cdot}}
\DeclareUnicodeCharacter{03BB}{\ensuremath{\lambda}}
\DeclareUnicodeCharacter{03B2}{\ensuremath{\beta}}
\DeclareUnicodeCharacter{2113}{\ensuremath{\ell}}
